\section*{Hoofdstuk \ref{ch:g6:pure-substances-and-mixtures} --- Zuivere stoffen en mengsels}

\begin{solution}{exo:g6:pure-substances-and-mixtures:1}
\omterm{def:g6:pure-substances-and-mixtures:pure-substance}{Zuivere stoffen}: gedestilleerd water, het blok zuiver ijzer, het
suikerkristal. \omterm{def:g6:pure-substances-and-mixtures:mixture}{Mengsels}: \omterm{def:g4:air-a-mixture-of-gases:air}{lucht}, zeewater, melk.
\end{solution}

\begin{solution}{exo:g6:pure-substances-and-mixtures:2}
Homogeen: zout water, kraanwater. Heterogeen: graniet, olie op water,
muesli.
\end{solution}

\begin{solution}{exo:g6:pure-substances-and-mixtures:3}
Eén bepaalde soort stof, met een eigen naam en eigen eigenschappen:
bijvoorbeeld water, zout, ethanol (ook suiker, zuurstof, ijzer).
\end{solution}

\begin{solution}{exo:g6:pure-substances-and-mixtures:4}
Een vast kookpunt wijst op een \omterm{def:g6:pure-substances-and-mixtures:pure-substance}{zuivere stof}; \qty{78}{\celsius} is het
kookpunt van ethanol.
\end{solution}

\begin{solution}{exo:g6:pure-substances-and-mixtures:5}
Het bevat veel \omterm{def:g6:pure-substances-and-mixtures:species}{chemische stoffen} (water, suikers, zuren, vitamines,
smaakstoffen): het is een \omterm{def:g6:pure-substances-and-mixtures:mixture}{mengsel}. ``Puur'' op het pak betekent alleen
dat er niets is toegevoegd.
\end{solution}

\begin{solution}{exo:g6:pure-substances-and-mixtures:6}
Nee: een \omterm{def:g6:pure-substances-and-mixtures:pure-substance}{zuivere stof} kookt bij een vaste temperatuur. Het is een
\omterm{def:g6:pure-substances-and-mixtures:mixture}{mengsel}, waarschijnlijk water waarin iets is opgelost, zoals zout water.
\end{solution}

\begin{solution}{exo:g6:pure-substances-and-mixtures:7}
Een filter kan het modderwater scheiden (de korreltjes blijven op het
papier) en het sap met vruchtvlees (het vruchtvlees blijft achter). Het
kan zout water niet scheiden (het zout is opgelost) en olie en water ook
niet (beide vloeistoffen gaan erdoorheen; je scheidt ze door de
bovenste laag af te gieten).
\end{solution}

\begin{solution}{exo:g6:pure-substances-and-mixtures:8}
$78.7 \div 10 = \qty{7.87}{g}$ voor \qty{1}{cm^3}: ijzer.
\end{solution}

\begin{solution}{exo:g6:pure-substances-and-mixtures:9}
$\frac{357}{478} \approx 0.747$: ongeveer \qty{75}{\percent}.
\end{solution}

\begin{solution}{exo:g6:pure-substances-and-mixtures:10}
$2 \times 35 = \qty{70}{g}$ zouten. Percentage: $\frac{35}{1000} =
\qty{3.5}{\percent}$.
\end{solution}

\begin{solution}{exo:g6:pure-substances-and-mixtures:11}
Meet hun kookpunten (elk moet tijdens het koken vast blijven, en de twee
moeten gelijk zijn), en weeg van elk hetzelfde volume (gelijke massa's).
Twee gelijke constanten, allebei vast, wijzen op dezelfde \omterm{def:g6:pure-substances-and-mixtures:pure-substance}{zuivere stof}.
\end{solution}

\begin{solution}{exo:g6:pure-substances-and-mixtures:12}
$\frac{10}{90 + 10} = \qty{10}{\percent}$ ethanol. Het is een \omterm{def:g6:pure-substances-and-mixtures:mixture}{mengsel}:
het heeft geen vast kookpunt en staat nergens in de tabel, die alleen
\omterm{def:g6:pure-substances-and-mixtures:pure-substance}{zuivere stoffen} vermeldt.
\end{solution}

\begin{solution}{pb:g6:pure-substances-and-mixtures:1}
\textbf{1.} Nee. Een \omterm{def:g6:pure-substances-and-mixtures:pure-substance}{zuivere stof} en een \omterm{def:g6:pure-substances-and-mixtures:homogeneous}{homogeen mengsel} kunnen er
precies hetzelfde uitzien: er moet iets gemeten worden.

\textbf{2.} Homogeen: het is helder en je ziet de delen niet.

\textbf{3.} Een stofconstante: een kookpunt dat vast blijft terwijl de
vloeistof kookt (of de massa van een bekend volume, vergeleken met een
tabel).

\textbf{4.} A en B, waarvan het kookpunt vast blijft, zijn \omterm{def:g6:pure-substances-and-mixtures:pure-substance}{zuivere
stoffen}. C is een \omterm{def:g6:pure-substances-and-mixtures:mixture}{mengsel}.

\textbf{5.} A kookt bij \qty{100}{\celsius}: water. B kookt bij
\qty{78}{\celsius}: ethanol.

\textbf{6.} A: $49.8 \div 50.0 = \qty{0.996}{g}$. B: $39.5 \div 50.0 =
\qty{0.790}{g}$. C: $51.2 \div 50.0 = \qty{1.024}{g}$.

\textbf{7.} Ja: ongeveer \qty{1.0}{g} voor water, \qty{0.79}{g} voor
ethanol.

\textbf{8.} Terwijl het water wegkookt, blijft het zout achter, dus de
vloeistof die overblijft bevat steeds meer zout; die kookt dan bij een
hogere temperatuur.

\textbf{9.} In vloeistof C is een vaste stof opgelost: het is een
oplossing, een \omterm{def:g6:pure-substances-and-mixtures:mixture}{mengsel}, zoals het koken al liet zien.

\textbf{10.} $100 \times 1.024 = \qty{102.4}{g}$; $\frac{3.5}{102.4}
\approx 0.034$: ongeveer \qty{3.4}{\percent} van de massa is zout.

\textbf{11.} $\frac{35}{1000} = \qty{3.5}{\percent}$: heel dicht bij
vraag 10.

\textbf{12.} \qty{3.5}{g} in \qty{100}{mL}, dus $3.5 \times 10 =
\qty{35}{g}$ zout per liter. Met zijn zoutgehalte en zijn kookgedrag is C
het zeewater; A is gedestilleerd water en B ethanol.
\end{solution}
